\chapter{The rationals}

\begin{notation}{The rationals}{}
  The rationals will be marked with a star. E.g., \(3\) would represent either an
  integer or a natural, but \(\Zrat{3}\) is a rational.
\end{notation}

\section{Construction of the rationals}

The rationals are constructed in a similar manner to the integers. Instead of the
\emph{a posteriori} knowledge of subtraction, we use the knowledge of division instead.

A rational is an ordered pair of integers:
\begin{equation}
  (a, b) = \frac{a}{b}
\end{equation}
where \(b \neq 0\). Two rationals \(\frac{a}{b}\) and \(\frac{c}{d}\) (\(b, d \neq 0\))
are equal \ifft
\begin{equation}
  ad = bc
\end{equation}
We then define the equivalence relation to be that way:

\begin{df}{Rational numbers}{rationals-definition}
  Define an equivalence relation
  \(\mathord{\sim} \subseteq \mathbb{Z} \times (\mathbb{Z} \smin \{0\})\) to be
  \begin{equation}
    (a, b) \sim (c, d) \mathiff ad = bc.
  \end{equation}
  A rational is then an equivalence class formed by \(\sim\). The set of all rationals
  is denoted \(\mathbb{Q}\).

  Let the sentence \(\Zrat{x} = (a, b)\) mean \((a, b) \in \Zrat{x}\).
\end{df}

\begin{prop}{The defining relation of the rationals is an equivalence relation}{}
  Let \((a, b), (c, d), (e, f) \in \mathbb{Z} \times (\mathbb{Z} \smin \{0\})\). The
  relation defined in \cref{df:rationals-definition} is an equivalence relation. I.e.,
  it is reflexive, symmetric, and transitive.
\end{prop}

\begin{proof}
  Since \(ab = ba\), \((a, b) \sim (a, b)\): it is reflexive.

  Assume that \((a, b) \sim (c, d)\), i.e., \(ad = bc\), or \(bc = ad\). Hence,
  \((c, d) \sim (a, b)\): it is symmetric.

  Now assume that \((a, b) \sim (c, d)\) and \((c, d) \sim (e, f)\). Then,
  \begin{equation}
    ad = bc \mathand cf = de.
  \end{equation}
  Multiplying the first equation by \(f\) and the second by \(b\) to get
  \begin{equation}
    adf = bcf \mathand bcf = bde,
  \end{equation}
  i.e., \(adf = bde\). By cancellation of multiplication on the integers, \(af = be\)
  which implies \((a, b) \sim (e, f)\): the relation is transitive.
\end{proof}

Arithmetic on the rationals is then defined as follows:

\begin{df}{Addition and multiplication on the rationals}{}
  Let \(\Zrat{x} = [(a, b)], \Zrat{y} = [(c, d)]\), then
  \begin{gather}
    \Zrat{x} + \Zrat{y} = [(ad + bc, bd)], \\
    \Zrat{x}\Zrat{y} = [(ac, bd)].
  \end{gather}
\end{df}

In which its justification for the equality symbol is again, another chorework.

\begin{prop}{Addition and multiplication on the rationals is well-defined}{}
  Let \(\Zrat{x} = [(a, b)]\), \(\Zrat{y} = [(c, d)]\) and \(\Zrat{z} = [(e, f)]\). If
  \((a, b) \sim (c, d)\), then
  \begin{equation}
    \Zrat{x} + \Zrat{z} = \Zrat{y} + \Zrat{z} \mathand \Zrat{x}\Zrat{z} = \Zrat{y}\Zrat{z}.
  \end{equation}
\end{prop}

\begin{proof}
  The relation \((a, b) \sim (c, d)\) means \(ad = bc\). Let's prove addition first.
  \begin{multline}
    \Zrat{x} + \Zrat{z} = [(af + be, bf)] \\
    \Zrat{y} + \Zrat{z} = [(cf + de, df)]
  \end{multline}
  For these two to be equal, we must have \((af + be, bf) \sim (cf + de, df)\) or
  \begin{equation}
    (af + be)(df) = (cf + de)(bf).
  \end{equation}
\end{proof}

\section{Ordering of the rationals}

\subsection{Absolute values}

\section{Max-min. Embedding the integers into the rationals}
